% Figure: planar magnetron sputtering source. Permanent magnets behind the
% target set up an arched magnetic field over the target face; the crossed
% radial E and the arched B trap electrons in a closed E-cross-B race track,
% intensifying ionisation just above the erosion groove. Ions are pulled to the
% target, sputter atoms off it, and the atoms deposit on the substrate above.
\begin{figure}[htbp]
\centering
\begin{tikzpicture}[>=Latex, scale=1.0]
  % target (cathode) slab
  \fill[enggray!25] (-3.2,0) rectangle (3.2,0.5);
  \node[enggray, font=\scriptsize, anchor=north east] at (3.2,0) {target (cathode)};
  % magnet bar cross-sections behind the target: outer poles S, centre pole N
  \fill[engblue!30] (-3.0,-0.7) rectangle (-2.2,0.0);
  \fill[engred!30]  (-0.4,-0.7) rectangle (0.4,0.0);
  \fill[engblue!30] (2.2,-0.7) rectangle (3.0,0.0);
  \node[font=\scriptsize] at (-2.6,-0.35) {S};
  \node[font=\scriptsize] at (0.0,-0.35) {N};
  \node[font=\scriptsize] at (2.6,-0.35) {S};
  % arched field lines from centre N to each outer S, above the target face
  \draw[engred, thick, ->] (-0.2,0.5) .. controls (-1.0,1.4) and (-1.8,1.4) .. (-2.6,0.5);
  \draw[engred, thick, ->] (0.2,0.5) .. controls (1.0,1.4) and (1.8,1.4) .. (2.6,0.5);
  % electron race-track marker: two crosses where field is most horizontal
  \node[engorange, font=\footnotesize] at (-1.4,1.15) {$\otimes$};
  \node[engorange, font=\footnotesize] at (1.4,1.15) {$\otimes$};
  \node[engorange, font=\scriptsize, anchor=south] at (0.0,1.5)
    {trapped-electron race track};
  % erosion grooves under the arches
  \draw[engblack, thick] (-1.7,0.5) .. controls (-1.4,0.28) and (-1.1,0.28) .. (-0.8,0.5);
  \draw[engblack, thick] (0.8,0.5) .. controls (1.1,0.28) and (1.4,0.28) .. (1.7,0.5);
  \node[engblack, font=\scriptsize, anchor=north] at (0.0,0.15) {erosion grooves};
  % sputtered atoms going up
  \draw[engblue, ->] (-1.25,0.6) -- (-1.0,2.4);
  \draw[engblue, ->] (1.25,0.6) -- (1.0,2.4);
  \draw[engblue, ->] (0.0,0.6) -- (0.0,2.4);
  % substrate at top
  \fill[englime!35] (-2.6,2.6) rectangle (2.6,3.0);
  \node[font=\scriptsize, anchor=south] at (0.0,3.0) {substrate (deposited film)};
  % incoming Ar+ ion pulled down to target
  \draw[engorange, ->] (2.0,2.2) -- (1.6,0.6);
  \node[engorange, font=\scriptsize, anchor=west] at (2.0,2.2) {Ar$^+$};
\end{tikzpicture}
\caption[Planar magnetron sputtering source]{A planar magnetron sputtering
source. Permanent magnets behind the target set up an arched magnetic field
(red) over the target face; the centre pole and the outer ring poles of
opposite sign make the field lines bow out and return, so above the erosion
groove the field runs nearly parallel to the target surface. Electrons emitted
from the negatively biased target are trapped there, gyrating on the field and
drifting azimuthally in the crossed radial electric and arched magnetic fields,
so they run in a closed race track (orange) that concentrates ionisation just
above the target. The dense plasma pulls argon ions into the target, where they
sputter atoms loose; the sputtered atoms travel up and deposit on the substrate
as a thin film. The magnetic trapping is what lets a magnetron run a dense
discharge at low pressure, which is why sputtered films are denser and cleaner
than those from an unmagnetised diode. The erosion grooves under the arches are
where the target wears fastest and set the target utilisation.}
\label{fig:vol04ch13:magnetron-sputter}
\end{figure}
